\documentclass[a4paper,12pt]{article}

\usepackage{a4wide}
\usepackage{amsfonts}
\usepackage{amsmath}
\usepackage{amssymb}
\usepackage{csquotes}
\usepackage{lipsum}
\usepackage{bm}
\usepackage{tabularx}
\usepackage{longtable}
\usepackage{booktabs}
\usepackage{array}
\usepackage{tikz}
\usepackage{float}
\usetikzlibrary{positioning, arrows.meta, shapes.geometric, fit}

\usepackage{graphicx}  % For including images
\usepackage{xcolor}  % For a colorfull presentation
\usepackage{listings}  % For presenting code 

\usepackage{algorithm}
\usepackage[noend]{algpseudocode}
\makeatletter
\def\BState{\State\hskip-\ALG@thistlm}
\makeatletter
\usepackage[normalem]{ulem}
\usepackage{array}
\usepackage{hyperref}
% \usepackage[hyphens]{url}

\hypersetup{
  colorlinks=true,
  linkcolor=blue,
  urlcolor=blue,
  citecolor=blue
}
\newcommand{\bluehref}[2]{\href{#1}{\textcolor{blue}{\uline{#2}}}}
\makeatletter
\newcommand*{\citelinktext}[2]{%
  \hyper@@link[cite]{}{cite.#1}{#2}}
\makeatother

\usepackage{booktabs} 
\usepackage{enumitem}
\usepackage{bbm}
\usepackage{tcolorbox}

\usepackage{minted}
\usepackage{subcaption}

\usepackage[font=small,labelfont=bf]{caption}

\usepackage{fullpage,amsmath,amsthm,amssymb}
\newtheorem*{solution}{Solution}
\newtheorem*{part1}{Part 1}
\newtheorem*{part2}{Part 2}

\newcommand{\Cov}{\mathrm{Cov}}
\newcommand{\Var}{\mathrm{Var}}
\newcommand{\rank}{\mathrm{rank}}
\newcommand{\Span}{\mathrm{span}}
\newcommand{\clip}{\mathrm{clip}}
\newcommand{\Dim}{\mathrm{dim}}

\newcolumntype{Y}{>{\raggedright\arraybackslash}X}
\newcolumntype{R}[1]{>{\raggedright\arraybackslash}p{#1}}

\tcbset{
  colback=gray!3,
  colframe=black!20,
  boxsep=1mm,
  left=1mm,
  right=1mm,
  top=1mm,
  bottom=1mm,
}

\usepackage{changepage}
\setlength{\parindent}{0pt}

\newenvironment{restoretext}%
    {
     \begin{adjustwidth}{}{\leftmargin}%
    }{\end{adjustwidth}
     }

% Definition of a style for code, matter of taste
\lstdefinestyle{mystyle}{
  language=Python,
  basicstyle=\ttfamily\footnotesize,
  backgroundcolor=\color[HTML]{F7F7F7},
  rulecolor=\color[HTML]{EEEEEE},
  identifierstyle=\color[HTML]{24292E},
  emphstyle=\color[HTML]{005CC5},
  keywordstyle=\color[HTML]{D73A49},
  commentstyle=\color[HTML]{6A737D},
  stringstyle=\color[HTML]{032F62},
  emph={@property,self,range,True,False},
  morekeywords={super,with,as,lambda},
  literate=%
    {+}{{{\color[HTML]{D73A49}+}}}1
    {-}{{{\color[HTML]{D73A49}-}}}1
    {*}{{{\color[HTML]{D73A49}*}}}1
    {/}{{{\color[HTML]{D73A49}/}}}1
    {=}{{{\color[HTML]{D73A49}=}}}1
    {/=}{{{\color[HTML]{D73A49}=}}}1,
  breakatwhitespace=false,
  breaklines=true,
  captionpos=b,
  keepspaces=true,
  numbers=none,
  showspaces=false,
  showstringspaces=false,
  showtabs=false,
  tabsize=4,
  frame=single,
}
\lstset{style=mystyle}

\begin{document}
\begin{titlepage}
    \centering
    \vspace*{1cm}
    
    \Huge
    \textbf{CSC490: Assignment 5}
    
    \vspace{0.5cm}
    
    \vspace{1.5cm}
    
    \textbf{Team Members}
    
    \vspace{0.8cm}
    
    \large
    \begin{tabular}{ll}
        \textbf{Name} & \textbf{Student ID} \\
        \midrule
        Aarya Prakash & 1008790257 \\
        Derek Huynh & 1008849052 \\
        Merrick Liu & 1008951920 \\
        Rhys Balevicius & 1000801974 \\
    \end{tabular}
    
    \vfill
    
    \large
    University of Toronto\\
    Department of Computer Science    
\end{titlepage}

\section{Context for Our Project: Lighthouse}

% ================================================================================
\section{Profiling Execution time}
% Profile 5 important functions in your codebase using https://docs.python.org/3/library/profile.html.
% Analyze your functions and find ways to improve them.

\textbf{Things to Profile: } 

\begin{enumerate}
    \item Search strategy (Hybrid, ...)
    \item Chunking strategy (Code, documentation)
    \item The MCP handler?
    \item Some of the temporal workflows? The could try to profile the indexing workflows one is pretty beefy
\end{enumerate}


% ================================================================================
\section{Showing Code Coverage}
% Using pytest-cov (https://github.com/pytest-dev/pytest-cov) or a similar library, show the code
% coverage of your project. Create a github workflow to generate your test coverage automatically on your
% readme and PRs. You can reference the example in pytest-cov or other open source projects

Quickly reference where the code coverage is.


% ================================================================================
\section{Test Coverage Breadth}
% Let’s make sure our code is well tested. You’ll recieve 0.20 marks for each 1% of code coverage up to 20
% marks based on part one.

Talk here about what we considered part of coverage.

% ================================================================================
\section{Test Coverage Depth}
% Choose your most complex module in your codebase. Create at least 10 unit tests for this module that
% test edge cases, failure modes and important use cases. Link the tests and edge cases, and explain why
% you chose these tests in comments of that module. We’ll try to find addtional edge cases that exist to
% cause an error! Make sure to write positive tests and negative tests (ones that throw errors or gracefully
% handle incorrect inputs)

Ingestion is th emost complex part.

% ================================================================================
\section{Two Memorable Bugs}
% Record a video walkthrough of two of the most memorable bugs you encountered during your project.
% For each bug, include:
% • What the bug was and how you discovered it
% • The debugging process you used to identify the root cause
% • How you fixed the bug
% • What you learned from the experience
% • The test cases you created to prevent this bug from reoccurring
% Link your video in your a5.pdf

Ingestion service. First one was what if multiple indices happen. Second one was incremental index deleting some stuff.


% ================================================================================
% \begin{thebibliography}{9}
% \bibitem{deepseekmath} Shao, Z., et al. (2024). DeepSeekMath: Pushing the Limits of Mathematical Reasoning in Open Language Models. arXiv preprint arXiv:2402.03300.
% \bibitem{nanochat-discussion}
% Karpathy, A. (2025). Nanochat Discussion \#1: Introducing nanochat: The best ChatGPT that \$100 can buy. GitHub Discussions. Available at: \url{https://github.com/karpathy/nanochat/discussions/1} (accessed March 11, 2026).

% \end{thebibliography}
\end{document}
